\documentclass[11pt,a4paper,fleqn]{article}
\usepackage{cancel}
\usepackage[left=1cm,right=1cm,top=1cm,bottom=1.8cm]{geometry}
\usepackage{graphicx}
\usepackage{titlesec}
\setlength{\mathindent}{0pt}
\titlespacing*{\subsection}{0pt}{1.2cm}{0.4cm}

\title{Double Pendulum Using Euler-Lagrange Equations and RK4}
\author{Assignment 1 MAS313}
\date{\today}

\begin{document}

\maketitle

\section{Introduction}

This project studies the motion of a double pendulum using the Euler-Lagrange
method to derive the equations of motion, followed by a fourth-order
Runge-Kutta method (RK4) to solve the resulting system numerically.

\section{Assumptions}

\begin{itemize}
    \item The masses are point masses.
    \item The rods are rigid and massless.
    \item The pendulum is ideal, with no loss from friction and no input from
    external forces.
\end{itemize}

\newpage

\section{Equations of Motion}

\begin{figure}[h]
    \centering
    \includegraphics[width=0.9\textwidth]{Screenshot 2026-09-16 183323.png}
\end{figure}

\subsection{Position}

\[
    x_1 = l_1 \sin(\theta_1)
\]
\[
    y_1 = -l_1 \cos(\theta_1)
\]
\[
    x_2 = l_1 \sin(\theta_1) + l_2 \sin(\theta_2)
\]
\[
    y_2 = -l_1 \cos(\theta_1) - l_2 \cos(\theta_2)
\]

\subsection{Velocity}

\[
    \dot{x}_1 = l_1 \dot{\theta}_1 \cos(\theta_1)
\]
\[
    \dot{y}_1 = l_1 \dot{\theta}_1 \sin(\theta_1)
\]
\[
    \dot{x}_2 = l_1 \dot{\theta}_1 \cos(\theta_1)
    + l_2 \dot{\theta}_2 \cos(\theta_2)
\]
\[
    \dot{y}_2 = l_1 \dot{\theta}_1 \sin(\theta_1)
    + l_2 \dot{\theta}_2 \sin(\theta_2)
\]

\newpage

\subsection{Potential Energy}

\[
    V = m_1 g y_1 + m_2 g y_2
\]

Substituting \(y_1\) and \(y_2\) into the expression gives

\[
    V = m_1 g \left(-l_1 \cos(\theta_1)\right)
    + m_2 g \left(-l_1 \cos(\theta_1) - l_2 \cos(\theta_2)\right)
\]

\[
    V = -m_1 g l_1 \cos(\theta_1)
    - m_2 g l_1 \cos(\theta_1)
    - m_2 g l_2 \cos(\theta_2)
\]

\[
    V = -(m_1 + m_2)g l_1 \cos(\theta_1)
    - m_2 g l_2 \cos(\theta_2)
\]

\subsection{Kinetic Energy}

\[
    T = \frac{1}{2}m_1\left(\dot{x}_1^2 + \dot{y}_1^2\right)
    + \frac{1}{2}m_2\left(\dot{x}_2^2 + \dot{y}_2^2\right)
\]

Substituting the velocities into the expression gives

\[
    T = \frac{1}{2}m_1 l_1^2 \dot{\theta}_1^2
    \cancelto{1}{\left(\cos^2(\theta_1) + \sin^2(\theta_1)\right)}
    + \frac{1}{2}m_2\left(\dot{x}_2^2 + \dot{y}_2^2\right)
\]

\[
    T = \frac{1}{2}m_1 l_1^2 \dot{\theta}_1^2
    + \frac{1}{2}m_2\left(
    \left(l_1 \dot{\theta}_1 \cos(\theta_1)
    + l_2 \dot{\theta}_2 \cos(\theta_2)\right)^2
    + \left(l_1 \dot{\theta}_1 \sin(\theta_1)
    + l_2 \dot{\theta}_2 \sin(\theta_2)\right)^2
    \right)
\]

\[
    \resizebox{\textwidth}{!}{$\displaystyle
    T = \frac{1}{2}m_1 l_1^2 \dot{\theta}_1^2
    + \frac{1}{2}m_2\left(
    l_1^2 \dot{\theta}_1^2
    \cancelto{1}{\left(\cos^2(\theta_1) + \sin^2(\theta_1)\right)}
    + l_2^2 \dot{\theta}_2^2
    \cancelto{1}{\left(\cos^2(\theta_2) + \sin^2(\theta_2)\right)}
    + 2l_1l_2\dot{\theta}_1\dot{\theta}_2
    \left(\cos(\theta_1)\cos(\theta_2)
    + \sin(\theta_1)\sin(\theta_2)\right)
    \right)$}
\]

\[
    T = \frac{1}{2}m_1 l_1^2 \dot{\theta}_1^2
    + \frac{1}{2}m_2\left(
    l_1^2 \dot{\theta}_1^2
    + l_2^2 \dot{\theta}_2^2
    + 2l_1l_2\dot{\theta}_1\dot{\theta}_2
    \cos(\theta_1 - \theta_2)
    \right)
\]

\newpage

\subsection{Lagrangian}

\[
    L = T - V
\]

Substituting \(T\) and \(V\) into the expression gives

\[
    L = \frac{1}{2}m_1 l_1^2 \dot{\theta}_1^2
    + \frac{1}{2}m_2\left(l_1^2 \dot{\theta}_1^2
    + l_2^2 \dot{\theta}_2^2
    + 2l_1l_2\dot{\theta}_1\dot{\theta}_2 \cos(\theta_1 - \theta_2)\right)
    - \left(-(m_1 + m_2)g l_1 \cos(\theta_1)
    - m_2 g l_2 \cos(\theta_2)\right)
\]

\[
    L = \frac{1}{2}m_1 l_1^2 \dot{\theta}_1^2
    + \frac{1}{2}m_2\left(l_1^2 \dot{\theta}_1^2
    + l_2^2 \dot{\theta}_2^2
    + 2l_1l_2\dot{\theta}_1\dot{\theta}_2 \cos(\theta_1 - \theta_2)\right)
    + (m_1 + m_2)g l_1 \cos(\theta_1)
    + m_2 g l_2 \cos(\theta_2)
\]

\subsection{Euler-Lagrange Equation}

\[
    \frac{d}{dt}\left(\frac{\partial L}{\partial \dot{q}_i}\right)
    - \frac{\partial L}{\partial q_i} = 0
\]

\subsubsection{Solving for the First Generalized Coordinate}

\[
    q_1 = \theta_1
\]

\[
    \frac{\partial L}{\partial \dot{\theta}_1}
    = (m_1 + m_2)l_1^2\dot{\theta}_1
    + m_2l_1l_2\dot{\theta}_2\cos(\theta_1 - \theta_2)
\]

\[
    \frac{d}{dt}\left(\frac{\partial L}{\partial \dot{\theta}_1}\right)
    = (m_1 + m_2)l_1^2\ddot{\theta}_1
    + m_2l_1l_2\left(
    \ddot{\theta}_2\cos(\theta_1 - \theta_2)
    - \dot{\theta}_2(\dot{\theta}_1 - \dot{\theta}_2)
    \sin(\theta_1 - \theta_2)
    \right)
\]

\[
    \frac{\partial L}{\partial \theta_1}
    = -m_2l_1l_2\dot{\theta}_1\dot{\theta}_2
    \sin(\theta_1 - \theta_2)
    - (m_1 + m_2)g l_1 \sin(\theta_1)
\]

\[
    \resizebox{\textwidth}{!}{$\displaystyle
    (m_1 + m_2)l_1^2\ddot{\theta}_1
    + m_2l_1l_2\left(
    \ddot{\theta}_2\cos(\theta_1 - \theta_2)
    - \dot{\theta}_2(\dot{\theta}_1 - \dot{\theta}_2)
    \sin(\theta_1 - \theta_2)
    \right)
    - \left(
    -m_2l_1l_2\dot{\theta}_1\dot{\theta}_2
    \sin(\theta_1 - \theta_2)
    - (m_1 + m_2)g l_1 \sin(\theta_1)
    \right) = 0$}
\]

\[
    \ddot{\theta}_1 =
    \frac{
    -m_2l_1l_2\ddot{\theta}_2\cos(\theta_1 - \theta_2)
    -m_2l_1l_2\dot{\theta}_2^2\sin(\theta_1 - \theta_2)
    -(m_1 + m_2)g l_1\sin(\theta_1)
    }{(m_1 + m_2)l_1^2}
\]

\subsubsection{Solving for the Second Generalized Coordinate}

\[
    q_2 = \theta_2
\]

\[
    \frac{\partial L}{\partial \dot{\theta}_2}
    = m_2l_2^2\dot{\theta}_2
    + m_2l_1l_2\dot{\theta}_1\cos(\theta_1 - \theta_2)
\]

\[
    \frac{d}{dt}\left(\frac{\partial L}{\partial \dot{\theta}_2}\right)
    = m_2l_2^2\ddot{\theta}_2
    + m_2l_1l_2\left(
    \ddot{\theta}_1\cos(\theta_1 - \theta_2)
    - \dot{\theta}_1(\dot{\theta}_1 - \dot{\theta}_2)
    \sin(\theta_1 - \theta_2)
    \right)
\]

\[
    \frac{\partial L}{\partial \theta_2}
    = m_2l_1l_2\dot{\theta}_1\dot{\theta}_2
    \sin(\theta_1 - \theta_2)
    - m_2g l_2 \sin(\theta_2)
\]

\[
    \resizebox{\textwidth}{!}{$\displaystyle
    m_2l_2^2\ddot{\theta}_2
    + m_2l_1l_2\left(
    \ddot{\theta}_1\cos(\theta_1 - \theta_2)
    - \dot{\theta}_1(\dot{\theta}_1 - \dot{\theta}_2)
    \sin(\theta_1 - \theta_2)
    \right)
    - \left(
    m_2l_1l_2\dot{\theta}_1\dot{\theta}_2
    \sin(\theta_1 - \theta_2)
    - m_2g l_2 \sin(\theta_2)
    \right) = 0$}
\]

\[
    \ddot{\theta}_2 =
    \frac{
    -l_1\ddot{\theta}_1\cos(\theta_1 - \theta_2)
    + l_1\dot{\theta}_1^2\sin(\theta_1 - \theta_2)
    - g\sin(\theta_2)
    }{l_2}
\]

\subsection{RK4 Setup}

The result is a system of two coupled differential equations, since the
acceleration of one angle depends on the acceleration of the other angle. Before
using RK4, the Python script rewrites the equations so that the accelerations
are explicit.

The expression for \(\ddot{\theta}_2\) is substituted into the expression for
\(\ddot{\theta}_1\), which gives an equation for \(\ddot{\theta}_1\) depending
only on the current angles and angular velocities. After \(\ddot{\theta}_1\) is
known, it is inserted back into the second equation to find
\(\ddot{\theta}_2\).

This gives the correct setup for RK4: at each RK4 step, and at each intermediate
stage, both accelerations are evaluated from the current state before advancing
the angles and angular velocities together.

\subsection{RK4 Method}

The numerical state is written as
\[
    \mathbf{y} =
    [\theta_1,\dot{\theta}_1,\theta_2,\dot{\theta}_2].
\]

The derivative of this state is
\[
    \frac{d\mathbf{y}}{dt} =
    [\dot{\theta}_1,\ddot{\theta}_1,\dot{\theta}_2,\ddot{\theta}_2].
\]

The RK4 method advances this state forward in time by evaluating the derivative
four times during each time step and combining the results.

\subsection{Energy Check}

Since the pendulum is assumed to be ideal, the total mechanical energy should be
constant. The kinetic energy, potential energy, and their sum are therefore
tracked during the simulation. Any change in the total energy is numerical
drift from the RK4 approximation.

\section{Conclusion}

The Euler-Lagrange method was used to derive the coupled equations of motion for
the double pendulum. These equations were then rewritten into a form suitable
for numerical simulation and solved using RK4. Energy was also tracked to check
how well the numerical solution conserved the ideal system energy.

\end{document}
